% =====================================================================
%  Math note template (compile with: pdflatex jacobian-vitushkin-fable.tex)
%  From Vitushkin's 2D example to Fable's 3D Jacobian counterexample
% =====================================================================
\documentclass[11pt]{article}

\usepackage[margin=1in]{geometry}
\usepackage{amsmath,amssymb,amsthm}
\usepackage{mathtools}
\usepackage{enumitem}
\usepackage[colorlinks=true,linkcolor=blue,urlcolor=blue,citecolor=blue]{hyperref}

% ---- theorem-like environments (handy for notes) --------------------
\theoremstyle{plain}
\newtheorem{proposition}{Proposition}
\newtheorem{lemma}{Lemma}
\theoremstyle{definition}
\newtheorem{remark}{Remark}

% ---- convenience macros ---------------------------------------------
\newcommand{\C}{\mathbb{C}}
\newcommand{\Jac}{\det DF}
\DeclareMathOperator{\Jdet}{J}

\title{From Vitushkin's Two-Dimensional Example\\
       to Fable's Three-Dimensional Jacobian Counterexample}
\author{}
\date{\today}

\begin{document}
\maketitle

\begin{abstract}
In July 2026 a polynomial map $\C^3\to\C^3$ with constant Jacobian
determinant $-2$ that is not injective was announced, disproving the
Jacobian Conjecture in every dimension $n\ge 3$. Its construction closely
mirrors an obscure 1999 Russian paper of Vitushkin. This note writes out,
step by step, a single-direction derivation of the three-dimensional
counterexample from Vitushkin's two-dimensional example, and shows that
degenerating the extra dimension collapses it exactly back to the
two-dimensional prototype.
\end{abstract}

% =====================================================================
\section*{Introduction}
% =====================================================================
In July 2026, Levent Alp\"oge, Akhil Mathew, and Anthropic's Claude Fable~5
produced a three-dimensional counterexample to the Jacobian Conjecture:
a polynomial map $\C^3\to\C^3$ whose Jacobian determinant is the nonzero
constant $-2$, yet which sends three distinct points to the same image.
Adjoining identity coordinates extends it to every $n\ge 3$, so the
conjecture is now false in all dimensions $\ge 3$; only the classical case
$n=2$ remains open.

This counterexample did not appear out of nowhere. Its construction closely
mirrors work of the Soviet/Russian mathematician Vitushkin in a 1999
Russian-language paper, which used $\sigma$-processes to build
two-dimensional rational maps with constant Jacobian that are globally
non-invertible, sketching the counterexample to within one step. Written in
Russian and long overlooked, it went unnoticed.

This note does one thing: it writes out, step by step, the derivation chain
from Vitushkin's two-dimensional framework to Fable's three-dimensional
counterexample --- from the two-dimensional fiber equation produced by the
$\sigma$-process, up one degree to three dimensions, through a birational
change of coordinates and a clearing of denominators, down to the three
polynomials --- and shows that degenerating the new dimension collapses it
exactly back to the two-dimensional prototype. The two examples are two
slices of a single construction, not unrelated objects.

% =====================================================================
\section*{Starting point: Vitushkin's two-dimensional example}
% =====================================================================
\[
  F = x^2y^6+2xy^2,\qquad G = xy^3+\frac{1}{y},\qquad \Jdet(F,G)=-2 .
\]
The obstruction is the $1/y$ in $G$: the Jacobian is a nonzero constant and
the map is globally non-invertible, but because of the denominator it is not
a polynomial, so it is not yet a counterexample to the Jacobian Conjecture.

% =====================================================================
\section*{Step 1: reduce the map to a single (quadratic) equation}
% =====================================================================
The essence of a counterexample is ``several inputs land on the same
output.'' So we ask the reverse question: fix an output; which inputs map to
it? These preimages are exactly the roots of a certain equation in an
auxiliary variable $T$ --- the \emph{fiber equation} (the equation for the
preimages of a fixed output). It is not arbitrary; it is the ledger of how
many preimages there are and where.

How does it arise from Vitushkin's original $F,G$? Through a
$\sigma$-\emph{process}, a coordinate substitution designed to ``look near
infinity'' (a blow-up type change) that spreads out the bad behavior
crowded at infinity and clears denominators along the way. Vitushkin builds
it from two pieces:
\[
  \Pi:\; x\mapsto \tfrac{x}{y},\; y\mapsto \tfrac1y;
  \qquad
  \Sigma_{xy}:\; x\mapsto xy,\; y\mapsto y .
\]
His combination $\Pi\cdot\Sigma_{xy}^4$ amounts to the single substitution
\[
  x\mapsto xy^3,\qquad y\mapsto \tfrac1y .
\]
The powers of $y$ it introduces cancel the high powers in $F$ and the $1/y$
in $G$:
\[
\begin{aligned}
  P &= F\!\left(xy^3,\tfrac1y\right)
     = (xy^3)^2\!\left(\tfrac1y\right)^6+2(xy^3)\!\left(\tfrac1y\right)^2
     = x^2+2xy,\\[2pt]
  Q &= G\!\left(xy^3,\tfrac1y\right)
     = (xy^3)\!\left(\tfrac1y\right)^3+\frac{1}{1/y}
     = x+y,
\end{aligned}
\]
with $\Jdet(P,Q)=2y$. The $\sigma$-process is not an arbitrary algebraic
manipulation but a choice of coordinate at infinity that exactly cancels the
denominators and high powers, turning the rational pair $F,G$ into the clean
polynomial pair $P,Q$.

Fixing the output $(P,Q)$ and solving for the preimages: from $Q=x+y$ we get
$y=Q-x$, and substituting into $P$,
\[
  P = x^2+2x(Q-x) = -x^2+2Qx
  \;\Longrightarrow\;
  x^2-2Qx+P = 0 .
\]
A quadratic, two roots, hence $2$-to-$1$. In the note's normalized output
variables $a,b$ (set $a=P$, $b=4Q$; the factor $4$ on the second coordinate
is the source of the ``off by a factor of $4$'' in the literature),
multiplying through by $-2$ gives the fiber equation
\begin{equation}
  -2\,T^2 + b\,T - 2a = 0 . \label{eq:fiber2d}
\end{equation}
Here $a,b$ are the outputs, $T$ is the ``label'' of a preimage (namely the
$x$ above), and each root corresponds to a preimage. Vitushkin is stuck
here: the denominator $1/y$ cannot be removed, and the degree of the
equation (two) is already saturated, with no higher term available to
``store'' the pole.

% =====================================================================
\section*{Step 2: lift by adding a cubic term, carried by a new coordinate $c$}
% =====================================================================
Add a cubic term $c\,T^3$ to the left of~(1), raising the degree to three.
Write this cubic polynomial as
\begin{equation}
  \Phi(T):=c\,T^3 - 2T^2 + bT - 2a,\qquad \Phi(T)=0 . \label{eq:fiber3d}
\end{equation}
\begin{itemize}[nosep]
  \item $c=0$: back to~(1), two-dimensional, $2$-to-$1$.
  \item $c\neq 0$: cubic, hence $3$-to-$1$.
\end{itemize}
The new coefficient $c$ is precisely the third output coordinate, carried by
a new variable $z$. This is the step ``2D $\to$ 3D.''

% =====================================================================
\section*{Step 3: tie the root to the pole}
% =====================================================================
Designate a root of~(2):
\[
  t = y + \frac{1}{x} = \frac{1+xy}{x}.
\]
Why this $t$? Because it is the reincarnation of the stubborn $1/y$ in
Vitushkin's example: the denominator $1/x$ is packed inside the root, carried
by the root rather than exposed in the map's formulas.

Now we must control \emph{branching}. The map is a three-sheeted cover (fix
an output; the three preimages are the three roots of $\Phi(T)=0$). Two
preimages collide --- the cover branches and the Jacobian degenerates ---
exactly when $\Phi$ has a double root, i.e.\ $\Phi(t)=0$ and $\Phi'(t)=0$.
Thus $\Phi'(t)$ is the ``where does it branch'' detector; differentiating in
$T$ is the standard test for a repeated root. Force $\Phi'(t)$ to be a pure
pole:
\begin{equation}
  \Phi'(t) = 3ct^2 - 4t + b = \frac{2}{x}. \label{eq:pole}
\end{equation}
Then for finite $x$ we have $\Phi'(t)=2/x\neq 0$, so branching can only occur
at $x=\infty$: the bad behavior is banished to infinity, and the finite part
is a clean three-sheeted cover (which is what makes a constant nonzero
Jacobian possible). This is the three-dimensional form of Vitushkin's
``push the exceptional curves to infinity.''

% =====================================================================
\section*{Step 4: solve back for the three source coordinates $(x,y,z)$}
% =====================================================================
The unknowns are the source coordinates $(x,y,z)$; the known quantities live
on the cubic-model side: the root $t$, the pole $r$, and the new output $c$.
This step joins the two coordinate systems --- a three-dimensional
birational change.

\medskip\noindent\textbf{$x$.} Set $r:=\dfrac2x$ (the pole on the right of
(3)), so $x=\dfrac2r$.

\medskip\noindent\textbf{$y$.} From the root $t=y+\dfrac1x$ and
$\dfrac1x=\dfrac r2$, we get $y=t-\dfrac r2$.

\medskip\noindent\textbf{$z$.} With no ready-made equation, write $z$ as an
ansatz with undetermined coefficients:
\[
  z = A\,r^2 + B\,tr + D\,c\,r^3 .
\]
The coefficients are pinned uniquely by two requirements (all verified
symbolically):
\begin{itemize}[nosep]
  \item \emph{Jacobian identically $-2$ fixes $D$.} For this shape one
        computes $\Jac=\dfrac{1}{4D}$ regardless of $A,B,D$: it is already
        constant (independent of $x,y,z$). Setting it to the target $-2$
        gives $D=-\dfrac18$.
  \item \emph{All three components polynomial fixes $A,B$.} With
        $D=-\dfrac18$, the component $c$ is automatically polynomial, but (see
        Step~5) $b$ carries a $\dfrac1x$ pole and $a$ a $\dfrac1{x^2}$ pole.
        Killing them gives $2A+B=1$ and $12A+8B=3$, whose unique solution is
        $A=\dfrac54,\ B=-\dfrac32$.
\end{itemize}
Hence
\[
  z=\frac54r^2-\frac32tr-\frac{c}{8}r^3 .
\]
The three relations
\[
  x=\frac2r,\qquad y=t-\frac r2,\qquad
  z=\frac54r^2-\frac32tr-\frac{c}{8}r^3
\]
are the three-dimensional birational change $(x,y,z)\leftrightarrow(t,r,c)$.

% =====================================================================
\section*{Step 5: substitute, clear denominators, obtain three polynomials}
% =====================================================================
Write the outputs $a,b,c$ in the source coordinates; each has a source.

\medskip\noindent\textbf{$b$ from (3).} Solving
$\Phi'(t)=3ct^2-4t+b=\dfrac2x$,
\[
  b = 4t+\frac2x-3ct^2 .
\]

\noindent\textbf{$2a$ from (2).} Solving $\Phi(t)=ct^3-2t^2+bt-2a=0$,
\[
  2a = ct^3-2t^2+bt .
\]

\noindent\textbf{$c$ by inverting the $z$-formula.} Solving
$z=\frac54r^2-\frac32tr-\frac c8r^3$ gives
$c=\dfrac{10}{r}-\dfrac{12t}{r^2}-\dfrac{8z}{r^3}$; with $r=\dfrac2x$ (so
$\dfrac1r=\dfrac x2$) this is $c=5x-3tx^2-x^3z$, and then $t=y+\dfrac1x$ gives
\[
  c = 5x-3\Big(y+\tfrac1x\Big)x^2-x^3z = 2x-3x^2y-x^3z .
\]

All three still carry $t=\dfrac{1+xy}{x}$ and $\dfrac1x$. Substituting into
$b,2a$ and clearing denominators, the $\dfrac1x$ cancels (this is where the
factor $(1+xy)$ appears), giving
\[
\begin{aligned}
  a &= (1+xy)^3z+y^2(1+xy)(4+3xy),\\
  b &= y+3x(1+xy)^2z+3xy^2(4+3xy),\\
  c &= 2x-3x^2y-x^3z .
\end{aligned}
\]
This is exactly Fable's three-dimensional example, with $\det=-2$.

% =====================================================================
\section*{The connection}
% =====================================================================
\[
  \underbrace{-2T^2+bT-2a=0}_{\text{Vitushkin, 2D, }2\text{-to-}1}
  \;\xrightarrow{\;+\,cT^3\ (\text{new coord }c\text{ carries }z)\;}\;
  \underbrace{cT^3-2T^2+bT-2a=0}_{\text{Fable, 3D, }3\text{-to-}1}
\]
Setting $c=0$ in Fable's cubic collapses it exactly to Vitushkin's
quadratic. Vitushkin's example is the $c=0$ slice of Fable's; Fable's example
is Vitushkin's equation lifted by a cubic term carried by $z$. This
single-direction derivation shows the two are homologous.

% =====================================================================
\section*{The special fiber: seeing the $3$-to-$1$}
% =====================================================================
At the output $(a,b,c)=\left(-\dfrac14,0,0\right)$ the cubic degenerates:
\[
  -2T^2+\frac12=0 \;\Longrightarrow\; T=\pm\frac12 .
\]
\begin{itemize}[nosep]
  \item Two finite roots give the preimages
        $\left(1,-\dfrac32,\dfrac{13}{2}\right)$ and
        $\left(-1,\dfrac32,\dfrac{13}{2}\right)$.
  \item The third root goes to $T=\infty$ (i.e.\ $x=0$), giving the preimage
        $\left(0,0,-\dfrac14\right)$.
\end{itemize}
Three preimages, matching $3$-to-$1$.

This point lies exactly on the two-dimensional seam. Its third output
coordinate is $c=0$, and, as noted above, $c=0$ collapses the
three-dimensional equation to the two-dimensional quadratic. Taking the first
two output coordinates $(a,b)=\left(-\dfrac14,0\right)$ down to two
dimensions, the fiber equation is the same $-2T^2+\dfrac12=0$ with the same
roots $T=\pm\dfrac12$. These are the two preimages of the two-dimensional
example at this output ($2$-to-$1$), corresponding in the $\sigma$-model
$P=x^2+2xy$, $Q=x+y$ to $(x,y)=\left(\dfrac12,-\dfrac12\right)$ and
$\left(-\dfrac12,\dfrac12\right)$. In three dimensions the same two finite
roots reappear, together with one root at infinity $T=\infty$ ($x=0$,
preimage $\left(0,0,-\dfrac14\right)$); this third sheet is what the extra
dimension adds. The special fiber is a live snapshot of ``3D $=$ 2D plus one
sheet at infinity'': the two finite sheets are already present in Vitushkin's
two-dimensional example, and the sheet at infinity is new to three
dimensions.

% =====================================================================
\section*{Discussion: Vitushkin's share is at least 60--70\%}
% =====================================================================
Laying out ``who contributed what,'' Vitushkin, the author of the 1999
Russian paper, takes the larger share, for concrete and checkable reasons.
\begin{enumerate}[label=\textbf{\arabic*.},nosep,leftmargin=*]
  \item \textbf{The whole machine is his.} The $\sigma$-process --- the
        coordinate change for computing Jacobians and building examples at
        infinity --- is the subject of his paper. The three-dimensional
        construction runs on it end to end.
  \item \textbf{The core local model is almost verbatim his.} Vitushkin's
        Example~1, after his own $\sigma$-process, gives $P=a^2+2ab$,
        $Q=a+b$ --- exactly the local model of the $c=0$ slice of Fable's
        example, off only by a factor of $4$ on the second coordinate. The
        engine producing ``constant Jacobian yet globally non-invertible'' is
        his, used almost unchanged.
  \item \textbf{The mechanism is his.} ``Constant nonzero Jacobian $+$
        globally many-to-one $+$ the bad behavior hidden in a pole'': this
        entire geometric mechanism he demonstrated in two dimensions. He built
        the bad map; only the pole remained uncleared.
  \item \textbf{He located the obstruction.} His Proposition~1 (a
        two-dimensional impossibility result) pins down exactly why two
        dimensions fail: the pole cannot be removed under these blow-ups.
        Knowing where the wall is tells you what to change --- add a
        dimension. He marked both the obstruction and the direction of the fix.
  \item \textbf{He was one step from the finish.} Method, a usable local
        model, the mechanism, the obstruction diagnosis --- he had all four;
        what remained was ``lift to three dimensions and balance the
        coefficients.''
\end{enumerate}

What Fable adds is the decisive last step: choosing a new variable $z$ to
absorb the pole, tuning the coefficients so that all three components become
polynomial with constant Jacobian, and actually turning the two-dimensional
``impossibility'' into a genuine three-dimensional counterexample. Anyone can
write the framework; producing the specific example with balanced
coefficients is the scarce step.

An honest delimitation: saying Vitushkin has $60$--$70\%$ refers to the
composition of the mathematical contribution (machine, model, mechanism, and
diagnosis are all his), not to a claim that Fable consciously copied him.
This note proves the two examples are homologous and mutually derivable (the
$c=0$ collapse, the term-by-term match of local models), which is a
verifiable fact; whether Fable actually retrieved and used the Russian paper
is a question of provenance that mathematics cannot settle, and is not
claimed here.

% =====================================================================
\section*{Conclusion}
% =====================================================================
Returning to the opening question: are Vitushkin's two-dimensional example
and Fable's three-dimensional counterexample related? The answer is yes, and
constructively so. The whole chain is: use Vitushkin's $\sigma$-process to
read off the two-dimensional fiber equation~(1); add a cubic term carried by
a new variable $z$ to lift it to the three-dimensional fiber equation~(2);
force the branch detector $\Phi'(t)$ to a pole~(3), banishing branching to
infinity; then solve back for the three-dimensional birational change and
clear denominators to obtain the three polynomials $a,b,c$ --- Fable's
three-dimensional counterexample.

The chain is one-way, explicit, and verifiable: (1), (2), (3), the
back-solving, and the clearing of denominators have all been checked
symbolically, and the final map satisfies $\det=-2$, $\Phi(t)=0$, and
$\Phi'(t)=2/x$. Setting $c=0$ collapses the three-dimensional fiber equation
exactly to Vitushkin's two-dimensional equation, whose local model matches
Vitushkin's 1999 Example~1 term by term (off by a factor of $4$ on the second
coordinate). The two examples are not lookalikes but the two-dimensional
slice and three-dimensional whole of a single construction.

Three honest caveats. First, ``why exactly this $\sigma$-chain'' is taken
from Vitushkin's paper and quoted, not re-derived. Second, the specific
birational change (especially the coefficients of $z$) is taken from the
community's reverse-engineering and may not be the original author's form,
with other equivalent forms possible. Third, this note proves mathematical
homology and mutual derivability, but does not assert Fable's actual
provenance. In sum, this three-dimensional counterexample pushes Vitushkin's
twenty-year-old two-dimensional near-miss forward by a decisive step; the
$n=2$ Jacobian Conjecture remains open, and remains the original and most
essential battleground of this two-dimensional machinery.

\end{document}
